\chapter{Discussion}
\label{chap:discussion}

This chapter interprets the Phase~2 results in relation to the research
questions and the literature.  The distinction between evidence,
interpretation, and implication is maintained throughout.  This distinction is
necessary because the campaign comprises one run per configuration
($N=1$), one graph seed, an eight-tick horizon, and 292 dead cells among
1,728 configured cells.  The reported quantities are therefore descriptive
properties of the present LLM simulation.  They are neither estimates of
population effects nor measurements of misinformation on Twitter.

\section{Synthesis of the research questions}
\label{sec:discussion-rqs}

\subsection{RQ1: What does the instrument comparison establish?}

The measurement result precedes all substantive comparisons. Dual-discrete scores exceeded continuous scores on every topology in both arms, while topology ranks aligned weakly. The campaign therefore shows that MI and MPR conclusions depend on the auditor regime. This finding answers RQ1: neither instrument can serve as an unqualified ground truth, and cross-mode pooling is not admissible.

\subsection{RQ3: How does network topology affect misinformation
propagation?}

\paragraph{Evidence.}
Topology changed the live-cell scores, the frequency of irreversible
\kstar{}, and the probability that a cascade remained scorable.  However, no
topology occupied the highest or lowest position consistently across identity
composition and scoring instrument.  In homogeneous continuous runs,
scale-free networks had the highest topology mean ($0.9205$), whereas
small-world networks had the lowest ($0.7215$).  In homogeneous
dual-discrete runs, echo chambers had the highest mean ($1.7984$), again with
small-world networks lowest ($1.5977$).  In heterogeneous runs, the
hierarchical topology had the highest mean under both instruments ($1.5242$
continuous and $2.5934$ dual-discrete).  These latter estimates must be read
together with the hierarchical dead-cell rates of $27.8\%$ and $22.2\%$,
respectively.  Polarised heterogeneous continuous runs likewise combined a
high live-cell mean ($1.4137$) with a $30.6\%$ dead-cell rate.

\paragraph{Interpretation.}
The results do not support a topology-invariant ranking.  Topology appears to
condition which messages are observed and scored, but it does so jointly with
persona placement, article content, and cascade survival.  In particular,
live-only means on topologies with many dead cells describe the surviving
cascades rather than the generator as a whole.  The scale-free versus random
comparison is primarily a hub and degree-sequence contrast; it is not a
controlled comparison of clustering coefficients.  Similarly, the
hierarchical generator represents authority-oriented transmission and should
not be interpreted as a model of cross-media dynamics.

\paragraph{Implication.}
Topology should be treated as a conditioning structure, not as a sufficient
cause of factual degradation.  Subsequent work should use repeated graph
seeds, report cascade survival as an outcome, and match graph properties when
isolating a specific mechanism such as clustering, centralisation, or
community separation.  On the present evidence, selecting a single
\enquote{most harmful} topology would overstate what the design identifies.

\subsection{RQ2: What is the role of identity composition?}

\paragraph{Evidence.}
When all eight generated topologies were collapsed within one instrument, the
heterogeneous arm exceeded the homogeneous arm: by $0.1549$ in continuous
\MI{} (H $0.8643$, He $1.0192$) and by $0.3986$ in dual-discrete \MI{}
(H $1.6814$, He $2.0800$).  The continuous direction was nevertheless split:
He exceeded H on four of eight topologies, whereas H exceeded He on the other
four.  The dual-discrete direction was He$>$H on all eight topologies.
Composition-specific results were substantially more differentiated than
these arm labels.  Homogeneous conspiracy-family cells reached $2.2194$
continuous and $3.1691$ dual-discrete, compared with $0.1118$ and $0.8687$
for homogeneous science/environment personas.  Among heterogeneous mixes,
\texttt{mix\_00}, containing four conspiracy personas, scored $1.7712$
continuous and $3.1321$ dual-discrete; conspiracy-free \texttt{mix\_02}
scored $0.2645$ and $1.1882$.

\paragraph{Interpretation.}
The collapsed H/He contrast is not an estimate of a general effect of
diversity.  The homogeneous arm contains both conspiracy-only and
scientist-only societies, while the heterogeneous arm contains mixes with
different conspiracy shares and different persona identities.  Averaging
within these broad labels changes the distribution of substantive roles.
Consequently, the observed He$>$H collapse is compatible with a strong
composition gradient: conspiracy-conditioned cells score high, whereas
science-conditioned cells score low.  It does not establish that
heterogeneity intensifies misinformation, just as the low value of
\texttt{mix\_02} does not establish that diversity prevents it.

\paragraph{Implication.}
Future comparisons should hold the proportion and identity of
conspiracy-conditioned seats constant while varying only within-network
diversity or placement.  Network means should be accompanied by
persona-conditioned \MPR{} and exposure-matched comparisons.  This would
separate a genuine interaction between heterogeneous neighbours from the
mechanical dilution or enrichment produced by changing the number of
high-scoring personas.

\subsection{Measurement implications across RQ2 and RQ3}

\paragraph{Evidence.}
Dual-discrete scores exceeded continuous scores on all eight topologies in
both arms.  The topology-mean difference was $0.8201$ for homogeneous and
$1.0450$ for heterogeneous runs.  Among cells live under both modes, the
median paired difference was $0.8125$ for H ($n=398$) and $0.9778$ for He
($n=198$).  Instrument choice also affected the incidence of irreversible
\kstar{}: the collapsed He rates were $32.1\%$ under dual-discrete scoring and
$7.6\%$ under continuous scoring.  Moreover, topology rankings agreed only
weakly across modes in the exploratory analysis (Spearman
$\rho=0.095$ for H and $\rho=0.167$ for He).

\paragraph{Interpretation.}
The two modes are not interchangeable measurements on a common numerical
scale.  The discrete instrument awards categorical penalties for omitted or
contradicted answers, whereas the continuous instrument resolves degradation
differently.  The systematic level difference and weak rank agreement show
that instrument choice can affect both absolute severity and comparative
claims about topology.  The continuous sidecar generated during a dual run is
useful for diagnosing disagreement, but it is not a third independent
\MPR{}.

\paragraph{Implication.}
Substantive conclusions should be reported separately for each instrument.
Averaging the two modes would create a quantity without a defined auditor.
The disagreement also motivates external validation: a preregistered sample
of rewrites should be assessed by human coders, and calibration should examine
whether each instrument preserves article- and persona-level ordering rather
than merely correlating at the event level.  Until such validation is
available, \kstar{} is instrument-relative.

\subsection{RQ4: Does the D-net analysis generalise the generated-topology
findings?}

\paragraph{Evidence.}
D-net used a 63-node custom graph reconstructed from hashtag co-occurrence,
not a retweet cascade.  Its homogeneous condition assigned
\texttt{conspiracy\_believer} to all 63 seats; its mixed condition retained a
conspiracy share of $28/63$.  On the identical edge set, homogeneous
conspiracy exceeded the mixed condition for both articles and both
instruments: continuous SCoPEx $1.7349$ versus $0.9983$, continuous
chemtrails--Gates $3.3378$ versus $2.5068$, dual-discrete SCoPEx $4.2537$
versus $2.6880$, and dual-discrete chemtrails--Gates $3.7303$ versus
$1.5297$.  Structural comparison was less favourable.  Breadth fell within
the specified $30\%$ band, whereas depth and structural virality did not;
DTFS was $0.2754$ and failed the $0.70$ validation threshold.  Tests comparing
one empirical graph with 16 simulated cascades were underpowered.

\paragraph{Interpretation.}
D-net provides a useful stress test of the composition interpretation because
the topology is held fixed while persona labels change.  It does not provide
out-of-sample validation of the eight-node generators, and it is not a
digital twin of Debnath et al.'s Twitter corpus.  The graph imports
co-occurrence structure and then simulates message propagation on that
structure; it neither reconstructs empirical retweet paths nor reproduces an
empirical misinformation score.  The D-net result therefore supports the
narrow statement that a conspiracy-only assignment scored above the retained
mixed assignment on this larger custom graph.

\paragraph{Implication.}
The qualitative composition gradient travels from the generated graphs to
D-net, but numerical effect sizes and cascade-shape claims do not.  Stronger
generalisation would require hydrated interaction data, common community and
centrality estimators, temporal information, multiple reconstructed cascades,
and repeated simulation seeds.  Simulated auditor \MI{} must remain separate
from Debnath et al.'s toxicity and network measures
\parencite{debnath2023conspiracy,ruths2014social}.

\section{Reconciling the collapsed and D-net orderings}
\label{sec:discussion-reconciliation}

\paragraph{Evidence.}
The generated-topology collapse gives He$>$H. In D-net, homogeneous
conspiracy scores exceed the mixed assignment. These comparisons use
different definitions of H.
Across the eight-node grid, H denotes twelve separate homogeneous persona
societies, including four conspiracy personas as well as climate-action,
environmental, journalist, and scientist personas.  D-net H denotes one
specific treatment: all 63 nodes are assigned the
\texttt{conspiracy\_believer} persona.  D-net He preserves a majority
non-conspiracy composition.  Restricting the generated-topology H cells to
the conspiracy family already reverses the broad collapse: conspiracy H is
well above overall He on both instruments.

\paragraph{Interpretation.}
There is therefore no empirical contradiction to resolve.  The label H
encodes within-network uniformity, but it does not encode which identity is
uniform.  Collapsing across homogeneous scientist and homogeneous conspiracy
societies answers a different question from comparing a conspiracy monoculture
with a mixed society.  D-net changes both diversity and the prevalence of the
high-scoring persona when moving from H to He.  Its ordering is consequently
consistent with composition, but cannot distinguish composition from a
causal benefit of heterogeneity.

\paragraph{Implication.}
H and He should not be used as portable substantive treatments without an
attached composition descriptor.  A defensible notation would distinguish,
for example, H-conspiracy, H-science, He-44\%-conspiracy, and
He-0\%-conspiracy.  A matched factorial design could then vary diversity while
holding conspiracy share, seed position, and topology fixed.  This
reformulation is more informative than asking whether H or He is
\enquote{higher} in the abstract.

\section{Relation to the award paper and adjacent literature}
\label{sec:discussion-literature}

\subsection{The CIKM/LASS study as prior method, not a benchmark result}

\paragraph{Evidence from the literature.}
\textcite{maurya2025simulating} introduced the auditor--node framework, MI and
MPR, and the severity taxonomy used by the present system.  Its published
evaluation concerned 30-hop linear branches, a different article collection,
a larger persona inventory, ten auditor questions, and a different model
configuration.  It reported that identity- and ideology-oriented homogeneous
branches could accelerate degradation and that experts could stabilise it;
it also reported propaganda-tier outcomes in approximately $85\%$ of
heterogeneous branch--domain pairs.  Thus, the award paper itself does not
establish a universal rule that heterogeneous societies must score below
homogeneous societies.

\paragraph{Interpretation of the present comparison.}
The present study transfers the instrument to climate and geoengineering
content and extends the setting from linear branches to generated and custom
graphs.  The non-zero, persona-sensitive scores indicate that the instrument
remains responsive in this domain.  The observed collapsed He$>$H ordering is
not a failed replication of a universal CIKM H$>$He hypothesis, because no
such universal hypothesis follows from the prior heterogeneous result.
Conversely, D-net H-conspiracy$>$He-mixed cannot be presented as a
replication: identity inventories, graph objects, hop budgets, auditor
questions, and model settings differ, and the present campaign has only one
replicate.

\paragraph{Implication.}
The defensible contribution is a boundary analysis of the prior method, not a
numerical reproduction of the award paper.  A direct conceptual replication
would require matched linear chains, an expert-inclusive heterogeneous arm,
the same climate articles in both arms, a sufficiently broad article set, and
node-type-conditioned outcomes.  The eight-hop graph campaign can motivate
that test but cannot substitute for it.

\subsection{Topology and diffusion literature}

\paragraph{Evidence from the literature.}
Classical diffusion models establish that network structure affects reach and
activation opportunities \parencite{kempe2003influence}.  Small-world and
preferential-attachment models formalise different structural properties
\parencite{watts1998collective,barabasi1999emergence}.  Empirical work further
shows that cascade depth and breadth are not captured by size alone
\parencite{goel2016virality}, and that false news can diffuse farther and
faster than true news in observational data \parencite{vosoughi2018false}.

\paragraph{Interpretation of the present comparison.}
The present results are compatible with the general proposition that
structure conditions diffusion, but they do not re-estimate these empirical
findings.  Here the state is mutable text, forwarding is mediated by an LLM,
and live-cell \MI{} is an auditor output.  Topology effects are also entangled
with cascade death and composition.  The data therefore add a
content-transformation perspective to network diffusion; they do not show
that one canonical graph class is intrinsically most susceptible to
misinformation.

\paragraph{Implication.}
Network controls in LLM societies should include both structural summaries
and content outcomes.  Matching degree, density, seed centrality, and
clustering would permit more specific tests than topology names alone.

\section{Pfeffer's firestorm theory: insights and boundaries}
\label{sec:discussion-pfeffer}

\paragraph{Evidence.}
\textcite{pfeffer2014firestorms} identify seven interrelated factors: speed
and volume of communication, binary choices, network clusters, unrestrained
information flow, lack of diversity, cross-media dynamics, and
network-triggered decision processes.  The present six-knob vocabulary is a
computational remapping rather than this original list.  The campaign varied
network structure and identity composition and compared article frames.  It
held surprise as a single-node drip, fixed the horizon at eight ticks, omitted
a cross-media layer, and used ternary actions---forward, reinterpret, and
drop---rather than binary platform actions.  D-net's empirical structure had
transitivity $0.315$, mean local clustering $0.515$, and mixed-label
homophily close to $0.84$, but these structures were pre-wired rather than
generated by the simulation.

\paragraph{Interpretation.}
Two Pfeffer insights receive limited support as useful design principles.
First, network clusters and identity structure cannot be reduced to a single
topology label: the high-scoring composition can be concentrated or diluted
within the same graph.  Second, lack of diversity is not adequately tested by
the binary H/He label unless identity prevalence is controlled.  The D-net
comparison illustrates both points particularly clearly.  By contrast, this
campaign does not test speed and volume as a shock manipulation,
unrestrained information flow at platform scale, binary-choice dynamics,
cross-media amplification, or the four stages of network-triggered adoption.
Nor does \kstar{} originate in Pfeffer's framework; it is a thesis-specific
auditor threshold.

\paragraph{Implication.}
Pfeffer's framework is most productive here as an audit of what the
simulation represents and omits.  It should not be presented as a
seven-factor confirmation.  A fuller test would add shock versus drip
seeding, a binary-action ablation, a legacy-media or exogenous-broadcast
layer, and explicit logging of repeated exposure and affirmation.  Such
extensions would move the framework from retrospective mapping towards
identified mechanism tests.

\section{Alternative explanations and validity of the interpretation}
\label{sec:discussion-validity}

Several alternative explanations remain plausible.  Persona prompts may
induce role compliance rather than persistent beliefs; high conspiracy scores
may therefore partly measure how faithfully the model follows an instruction
to use conspiratorial language.  The action policy gives reinterpretation a
substantial role, which creates opportunities for factual mutation by design.
The same model family participates in generation and automated assessment,
and the five author-defined questions constrain what counts as degradation.
Article differences combine topic, wording, affective charge, and prior fit
with persona prompts, so they are not a clean manipulation of valence.
Finally, excluding dead cells from means avoids coding them as factual zeros,
but produces survivor-conditioned estimates.

These limitations do not make the campaign uninformative.  They narrow the
appropriate claim.  The experiment demonstrates that, under the documented
prompts and graph protocols, auditor scores vary with persona composition,
topology, and scoring mode.  It does not establish human susceptibility,
platform-level firestorm dynamics, or causal effects in an empirical
population.  Human evaluation, model ablations, repeated seeds, and a
fail-closed auditor are therefore validation requirements rather than
optional refinements \parencite{zheng2023judge,park2023generative}.

\section{Overall implications}
\label{sec:discussion-implications}

For simulation methodology, the principal implication is that treatment
labels must preserve their substantive content.  \enquote{Homogeneous} and
\enquote{heterogeneous} are insufficient without persona shares, placement,
and node-conditioned outcomes.  Likewise, topology names should be
accompanied by realised graph statistics and cascade-survival rates.

For measurement, discrete and continuous \MI{} should be treated as separate
operationalisations.  Agreement in a direction across both instruments is
stronger descriptive evidence than a result present under only one, but
neither becomes ground truth through convergence.  Divergence is itself a
result about the sensitivity of the conclusion to the auditor.

For empirical generalisation, D-net is an intermediate case: larger and
structurally grounded in a hashtag graph, yet still neither a retweet cascade
nor a temporal reconstruction.  Its contribution is to show that the
composition interpretation survives a change in graph size and source under
the same simulator.  It does not validate the simulator against Twitter.

Taken together, the research questions yield a restrained conclusion.
Topology conditions propagation but has no stable winner; identity
composition is more informative than the H/He label; severity and \kstar{}
are instrument-dependent; and D-net supports only qualitative, not numerical,
generalisation.  Pfeffer's framework explains why these dimensions should be
considered jointly, while also making visible the mechanisms that the current
campaign leaves untested.
